\chapter{Asymmetric Cryptography}
\label{ch:5.4}
\chaptertopic{Topic 5.4}

\chapterquote{We stand today on the brink of a revolution in cryptography.}{Whitfield Diffie and Martin Hellman, ``New Directions in Cryptography'' (1976)}

\section{Motivation}

Symmetric encryption in Chapter~\ref{ch:5.3} has an uncomfortable bootstrap problem: both parties
need the same secret key before they can talk securely. So how do they share it the first time, over
a network anyone might be listening to? This chapter's answer is asymmetric cryptography: two
mathematically linked keys, one of which can be handed to literally anyone without weakening
security at all. The chapter builds from the key pair itself, through the specific algorithm (RSA)
that makes it work, to digital signatures. It closes with the infrastructure that lets you trust a
stranger's public key really belongs to them in the first place.

\begin{apalignment}{Unit 5: Securing Applications and Data}{Topic 5.4 -- Asymmetric Cryptography}
\begin{apobjectives}
\apobjective{5.4.A}{Mitigate Risk}{Use a public/private key pair correctly for the operation actually intended -- encrypting to a recipient's public key for confidentiality, or signing with one's own private key for authenticity -- recognizing that the two operations are not reversible readings of one another.}
\apobjective{5.4.B}{Analyze Risk}{Compare key sizes across algorithm families on a common footing (security strength in bits, per NIST SP 800-57) rather than assuming equal key length means equal security, and reason about what a given key size still protects against today's best-known attacks.}
\apobjective{5.4.C}{Mitigate Risk}{Establish trust in a stranger's public key through a certificate and a chain of trust anchored in a pre-installed root, and evaluate whether a presented certificate should actually be trusted.}
\end{apobjectives}
\end{apalignment}

\section{Public/Private Key Pairs}

\begin{definition}{Asymmetric cryptography}
A cryptographic system built on two mathematically linked keys -- a \textbf{public key} that can be
freely distributed, even posted publicly, and a \textbf{private key} kept secret by its owner
forever -- where what one key does, only its partner can undo.
\end{definition}

Key generation starts from a hard mathematical problem: easy to compute in one direction,
computationally infeasible to reverse without special knowledge. For RSA, that problem is integer
factorization -- multiplying two large primes is fast, but recovering them from only their product is
not, for large enough primes. Two large primes $p$ and $q$ are chosen; their product $n = p \times q$
becomes part of both keys. A public exponent $e$ is chosen (nearly always 65537 in practice).
A private exponent $d$ is then derived such that $e \times d \equiv 1 \pmod{\phi(n)}$, where
$\phi(n) = (p-1)(q-1)$ -- a value that can only be computed efficiently if you already know $p$ and
$q$. The public key is $(e, n)$; the private key is $(d, n)$.

The keys are asymmetric in \emph{function}, not merely in secrecy, and this produces two entirely
separate use cases. \textbf{Encrypt with the public key; decrypt with the private key}: if Alice
wants to send Bob a secret, she encrypts with Bob's public key. Only Bob's private key can
reverse it -- not even Alice, who did the encrypting, can decrypt her own message afterward. This is
confidentiality. \textbf{Sign with the private key; verify with the public key}: if Bob wants to
prove he authored something, he signs with his own private key. Anyone holding his public key
can then verify it, knowing only Bob could have produced a signature that checks out. This is
authenticity, covered fully later in this chapter.

\begin{example}[The one-way trapdoor]
A useful mental model is a trapdoor in a floor: easy to fall through in one direction, impossible to
climb back through without the special mechanism. Multiplying $p$ and $q$ is the easy fall; factoring
$n$ back into $p$ and $q$ is the locked climb. The private key is the trapdoor mechanism that
makes the reverse direction possible for its holder alone. Other asymmetric systems use a different
one-way problem entirely -- discrete logarithms for Diffie-Hellman, elliptic-curve discrete logarithms
for ECC -- but the trapdoor structure recurs universally.
\end{example}

\begin{lstlisting}[language=Python, caption={Deriving and validating RSA key parameters from two primes}]
from math import gcd

def derive_key_params(p, q, e):
    """Returns (n, phi_n) and validates that e is a legal public exponent."""
    n = p * q
    phi_n = (p - 1) * (q - 1)
    if not (1 < e < phi_n) or gcd(e, phi_n) != 1:
        raise ValueError("e is not a valid public exponent for this n")
    return n, phi_n

def is_valid_key_pair(e, d, phi_n):
    return (e * d) % phi_n == 1
\end{lstlisting}

Security scales with key size, since larger keys make the underlying factoring problem harder. But
key length is not comparable across algorithm families: a 256-bit AES key, a 256-bit ECC key, and a
256-bit RSA modulus describe three completely different security levels, because each family faces a
different best-known attack. The unit that \emph{is} comparable is \textbf{security strength},
measured in bits, where a scheme with $s$ bits of strength requires about $2^s$ operations to break.
Table~2 of NIST SP 800-57 Part~1 Rev.~5\srccite{5.4}{1} gives the mapping between families:

\begin{center}
\begin{tabular}{cccc}
\toprule
\textbf{Strength (bits)} & \textbf{AES} & \textbf{RSA modulus} & \textbf{ECC key} \\
\midrule
112 & 3-key TDEA / AES-112-equiv & 2048-bit & 224--255 bit \\
128 & AES-128 & 3072-bit & 256--383 bit \\
192 & AES-192 & 7680-bit & 384--511 bit \\
256 & AES-256 & 15360-bit & $\geq$512 bit \\
\bottomrule
\end{tabular}
\end{center}

The RSA column explodes while the other two grow roughly linearly because of the shape of the best
attack against each. Against a well-chosen symmetric cipher or elliptic curve, the best known attack
is generic -- exhaustive search for AES, Pollard's rho for ECC, which costs about the square root of
the group order. So one more bit of key buys one more bit of strength for AES, and roughly half a bit
for ECC. RSA instead faces the General Number Field Sieve, whose running time is
\emph{sub-exponential} in the modulus size. So the next 16 bits of strength cost not a few more bits
of modulus but thousands: 2048 to 3072 bits for the step from 112 to 128 bits of strength, and 3072
to 15360 bits for the step from 128 to 256. RSA private-key operations also cost roughly the cube
of the modulus size, so a 15360-bit RSA key is not merely large to transmit but unusable at scale,
while a 512-bit elliptic curve remains cheap. That cost curve is why modern TLS and SSH deployments
default to ECDSA P-256 or Ed25519 rather than to ever-larger RSA keys. It is also why an RSA-1024
key, offering only about 80 bits of strength, sits below the security floor NIST withdrew support
for in 2013.\srccite{5.4}{2} None of
this matters, however, unless you actually receive the \emph{genuine} public key belonging to whoever
you want to talk to, rather than one an attacker substituted -- the trust problem this chapter closes
with in its final section.

\begin{misconception}
\textbf{``If I encrypt with the private key, it is more secure.''} Encrypting with the private key
provides no confidentiality at all, since anyone holding the matching public key can decrypt it.
Signing with the private key is for authenticity, a completely different guarantee from secrecy.
\end{misconception}

\section{The RSA Algorithm}

Everything RSA does reduces to one operation, applied twice: \textbf{modular exponentiation}.
Encryption computes $c = m^e \bmod n$; decryption computes $m = c^d \bmod n$. A theorem from number
theory (built on Euler's theorem) guarantees $(m^e)^d \bmod n = m$ whenever $d$ was derived correctly
relative to $e$, $p$, and $q$. Encrypting and then decrypting always returns the original message,
but only for whoever holds the matching $d$.

\begin{example}[A complete worked example, with nothing simplified away]
\textbf{Key generation.} Choose two primes $p = 61$ and $q = 53$. The modulus is
$n = pq = 3233$. Euler's totient is $\phi(n) = (p-1)(q-1) = 60 \times 52 = 3120$. Choose a public
exponent coprime to $\phi(n)$: take $e = 17$, since $\gcd(17, 3120) = 1$. The private exponent $d$ is
the modular inverse of $e$ modulo $\phi(n)$, i.e.\ the value satisfying $ed \equiv 1 \pmod{3120}$.
The extended Euclidean algorithm gives $d = 2753$, which checks: $17 \times 2753 = 46801 = 15 \times
3120 + 1$. The public key is $(e,n) = (17, 3233)$; the private key is $(d,n) = (2753, 3233)$. The
primes $p$ and $q$ themselves must now be destroyed or protected exactly as carefully as $d$, since
either one immediately yields $d$.

\textbf{Encryption.} To send $m = 65$, a sender computes $c = m^e \bmod n$ using the recipient's
public key. Doing this by square-and-multiply, so every intermediate value can be checked by hand:
\[
65^1 \bmod 3233 = 65, \quad 65^2 \bmod 3233 = 992, \quad 65^4 \bmod 3233 = 1232,
\]
\[
65^8 \bmod 3233 = 1547, \quad 65^{16} \bmod 3233 = 789.
\]
Since $17 = 16 + 1$, $65^{17} = 65^{16} \times 65^1 \equiv 789 \times 65 = 51285 \equiv 2790
\pmod{3233}$. The ciphertext is $c = 2790$.

\textbf{Decryption.} The recipient recovers $m = c^d \bmod n = 2790^{2753} \bmod 3233 = 65$ using
their private key -- an exponent large enough that it is checked by machine rather than by hand.
Notice that running the same pair of operations in the opposite order, private key first and public
key second, is the skeleton of a signature rather than of an encryption. The direction of the two
keys, not the arithmetic itself, is what distinguishes the two use cases.
\end{example}

\begin{lstlisting}[language=Python, caption={RSA encryption and decryption via modular exponentiation}]
def rsa_encrypt(m, e, n):
    return pow(m, e, n)

def rsa_decrypt(c, d, n):
    return pow(c, d, n)
\end{lstlisting}

Real implementations almost universally use $e = 65537$ ($2^{16}+1$) rather than a small arbitrary
prime, for three reasons. It is prime, so $\gcd(e, \phi(n)) = 1$ is essentially guaranteed for random
large $p, q$. Its binary form (a single 1 bit, fifteen 0 bits, a final 1 bit) makes repeated-squaring
exponentiation fast. And it is large enough to defeat low-exponent attacks that break RSA when $e$ is
very small (such as $e=3$) and the same message is sent unpadded to multiple recipients. Why this
resists reversal without $d$ traces directly back to Chapter~\ref{ch:5.4}'s first section: computing
$d$ from $(e, n)$ alone requires $\phi(n)$, which requires factoring $n$ into $p$ and $q$. That is
trivial for a toy $n=33$, but computationally infeasible for a real 2048-bit modulus with no known
classical algorithm. The asymmetry that protects RSA lives entirely in the difficulty of factoring,
never in the exponentiation itself, which is fast in both directions once the correct exponent is
known.

\begin{misconception}
\textbf{``Small numbers in the worked example mean RSA is easy to break.''} The arithmetic is
identical at any size; security comes from using primes hundreds of digits long, where factoring $n$
becomes computationally infeasible -- not from the operations themselves becoming more complex.
\end{misconception}

\section{Digital Signatures}

A digital signature is the asymmetric analogue of the HMAC from Chapter~\ref{ch:5.3}. Where HMAC
relies on one shared secret both parties possess, a signature uses a key pair: only the signer holds
the private key, and anyone at all can hold the public key. That structural difference unlocks a
guarantee HMAC cannot offer.

Signing the entire message directly with RSA's private-key operation is both slow (RSA exponentiation
cost grows with input size) and structurally risky. An attacker who obtains signatures on certain
chosen messages can sometimes combine them algebraically to forge a signature on a message never
actually signed. Real schemes fix both problems by never signing the message itself:

\begin{lstlisting}[language=Python, caption={Signing and verifying via a hashed digest}]
import hashlib

def sign(message, d, n):
    digest = int(hashlib.sha256(message.encode()).hexdigest(), 16) % n
    return pow(digest, d, n)

def verify(message, signature, e, n):
    expected_digest = int(hashlib.sha256(message.encode()).hexdigest(), 16) % n
    recovered_digest = pow(signature, e, n)
    return recovered_digest == expected_digest
\end{lstlisting}

The signer computes a fixed-size digest of the message with a cryptographic hash, and signs
\emph{that digest} rather than the message. A verifier recomputes the hash independently, recovers
the digest the signature actually encodes, and checks the two match. Hashing first keeps the
signature a fixed, small size regardless of message length, and runs the expensive RSA step only
once. It also inherits the avalanche effect from Chapter~\ref{ch:5.3}: a single changed bit in the
message produces a completely different hash, so any tampering is caught the instant verification
is attempted.

A signature that verifies successfully gives three distinct guarantees at once. \textbf{Authenticity}
-- the signature could only have been produced by whoever holds the matching private key.
\textbf{Integrity} -- because the signature covers a hash of the message, any post-signing
modification changes that hash and breaks verification immediately. \textbf{Non-repudiation} -- the
signer cannot credibly deny having signed, since nobody else possesses their private key.

\begin{example}[Why signatures beat HMAC for this specific guarantee]
HMAC gives integrity and authenticity between exactly two parties who share a secret key, but no
non-repudiation at all. If Alice and Bob share an HMAC key and Alice receives a validly tagged
message, she knows it came from ``Bob or herself,'' since either party holding the shared secret
could have produced that exact tag. There is no way to prove to a third party which of the two
actually created it. A digital signature removes this ambiguity structurally, because signing and
verifying use \emph{different} keys: only the signer can produce a valid signature, but anyone can
verify one. That also lets a single public key scale to an unlimited number of verifiers, unlike
HMAC, where every new verifier would need its own shared secret -- and could then forge tags for
everyone else.
\end{example}

\begin{misconception}
\textbf{``A digital signature encrypts the message for confidentiality.''} It provides no secrecy at
all; the message is frequently sent in the clear alongside its signature. Signing proves authenticity
and integrity, never confidentiality -- that requires separately encrypting with the recipient's
public key.
\end{misconception}

\section{PKI and Digital Certificates}

The previous section left a bootstrap problem deliberately unsolved: asymmetric cryptography lets
you hand out your public key freely, but only helps if the recipient can be sure that key genuinely
belongs to who it claims. If an attacker slips you a different public key while claiming to be your
bank, every guarantee from this chapter collapses at once -- you would encrypt your password with the
attacker's key, believing it protected.

\begin{definition}{Digital certificate}
A small, signed document binding a public key to an identity, containing at minimum a
\textbf{subject} (the identity vouched for), the subject's public key, an \textbf{issuer} (who is
vouching), a validity period, and a digital signature over all of those fields, produced with the
issuer's own private key.
\end{definition}

That signature is exactly the sign/verify operation from the previous section: the issuer hashes the
certificate's fields and signs the digest. Anyone holding the issuer's public key can then verify
that the certificate has not been altered since the issuer vouched for it. The issuer here is a
\textbf{Certificate Authority (CA)}, but this only pushes the trust question up one level: why trust
the CA's own public key?

\begin{definition}{Chain of trust}
The sequence of signed certificates from a leaf certificate up through intermediate CAs to a
\textbf{root CA}, whose self-signed certificate is pre-installed in an operating system or browser's
\textbf{trust store}, verified out-of-band by the OS vendor rather than by the individual user.
\end{definition}

A root CA rarely signs end-entity certificates directly -- its private key is kept extremely well
protected, often entirely offline -- and instead signs \textbf{intermediate CA} certificates,
delegating signing authority downward. The intermediate CA then signs the actual leaf certificate
presented to a browser. Validating this chain means verifying each certificate's signature using the
\emph{next} certificate's public key -- the leaf against the intermediate's key, the intermediate
against the root's. The validator then confirms the chain terminates at a root already present in
the trust store.

\begin{figure}[h]
\centering
\begin{tikzpicture}[
  cert/.style={draw=kickstartermaroon, thick, rounded corners, minimum width=3.2cm, minimum height=1cm, align=center, font=\scriptsize},
]
\node[cert, fill=kickstartermaroon!25] (root) at (0,2.4) {Root CA\\self-signed, in OS/browser\\trust store};
\node[cert, fill=kickstartercream] (inter) at (0,1.2) {Intermediate CA\\signed by root's\\private key};
\node[cert] (leaf) at (0,0) {Leaf certificate\\(the website)\\signed by intermediate's\\private key};
\draw[-{Latex[length=2mm]}, kickstartermaroon, thick] (inter) -- (root) node[midway, right, font=\tiny] {verify against};
\draw[-{Latex[length=2mm]}, kickstartermaroon, thick] (leaf) -- (inter) node[midway, right, font=\tiny] {verify against};
\node[font=\scriptsize\sffamily, gray, align=left] at (4.6,1.2) {each arrow is a signature\\checked with the\\\emph{next} cert's public key};
\end{tikzpicture}
\caption{The chain of trust: validation walks from the leaf certificate up to a root already
present in the trust store, verifying each signature with the next certificate's public key.}
\label{fig:5.4-chain-of-trust}
\end{figure}

\begin{lstlisting}[language=Python, caption={Walking a certificate chain to verify trust}]
def validate_chain(chain, trusted_roots, verify_signature):
    """chain: list of cert dicts ordered leaf-first, each like
       {"subject": str, "issuer": str, "public_key": ..., "signature": ...,
       "signed_fields": ..., "not_before": int, "not_after": int}.
       trusted_roots: set of root subject names.
       verify_signature(cert, issuer_public_key) -> bool.
       Returns (True, None) or (False, reason)."""
    for i, cert in enumerate(chain):
        if i + 1 < len(chain):
            issuer_cert = chain[i + 1]
            if cert["issuer"] != issuer_cert["subject"]:
                return False, "broken-chain-link"
            if not verify_signature(cert, issuer_cert["public_key"]):
                return False, "invalid-signature"
        else:
            if cert["subject"] not in trusted_roots:
                return False, "untrusted-root"
            if cert["subject"] != cert["issuer"]:
                return False, "root-not-self-signed"
    return True, None
\end{lstlisting}

Beyond the signature chain, a client checks each certificate's \textbf{validity window} -- one
presented before \texttt{not\_before} or after \texttt{not\_after} is rejected, which is why an
expired certificate throws a browser warning. A client also checks its \textbf{revocation status},
since a certificate's stated expiration might be years away while the situation that made it
trustworthy can change immediately, most commonly because its private key was stolen. A CA cannot
un-sign an already-issued certificate, so it instead publishes revocation through a \textbf{Certificate
Revocation List (CRL)}, a periodically updated, CA-signed list of revoked serial numbers. Or it
publishes through \textbf{OCSP}, where a client asks a CA-operated responder in real time whether one
specific certificate is still good and receives a signed answer without downloading an entire list.

\begin{misconception}
\textbf{``Revocation and expiration are the same thing.''} Expiration is a scheduled, built-in end
to a certificate's validity window. Revocation is an unscheduled, active withdrawal of trust before
that date, typically because the private key was compromised in the meantime.
\end{misconception}

\section*{Review Questions}
\begin{reviewquestions}
\item Alice encrypts a message with Bob's public key and sends it. Explain why Alice herself cannot
decrypt the ciphertext afterward, and explain what property of the key pair makes this true.
\item Using the worked RSA example in this chapter ($p=61$, $q=53$, $e=17$, $d=2753$, $n=3233$),
explain in your own words why encrypting $m=65$ and then decrypting the result returns exactly $65$,
referencing the relationship between $e$, $d$, and $\phi(n)$.
\item Explain why real signature schemes hash a message before signing it, rather than signing the
raw message directly with RSA, citing both a performance reason and a structural-security reason from
this chapter.
\item Two colleagues, Alice and Bob, share an HMAC key. Alice receives a validly tagged message and
cannot prove to a third party whether Bob or she herself produced it. Explain why switching to
digital signatures resolves this specific ambiguity.
\item A browser encounters a self-signed certificate for a public website and rejects it, while the
same browser accepts a different self-signed certificate manually installed on a private internal
network. Explain why the identical cryptographic construction produces two different trust outcomes.
\item Using the chain-of-trust diagram, explain what specifically breaks validation if an
intermediate CA's certificate is presented but the leaf certificate's \texttt{issuer} field does
not match the intermediate's \texttt{subject} field, even if every individual signature in the
chain is otherwise mathematically valid.
\item A root CA's private key is compromised, but the root CA itself has not yet issued a
revocation or informed browser vendors. Explain why the security strength table earlier in this
chapter (comparing AES, RSA, and ECC key sizes) is irrelevant to assessing the danger here, and
identify what actually determines the scope of damage.
\end{reviewquestions}

\begin{answerkey}
\item Alice cannot decrypt her own ciphertext because encryption with Bob's public key can only be
reversed by Bob's matching private key, which only Bob holds -- not even the party that performed the
encryption. This follows from the defining property of the key pair: what one key does, only its
mathematically linked partner can undo, and Alice possesses neither Bob's private key nor any
shortcut around needing it.
\answerlabel{5.4.A}
\item Because $d$ was derived so that $ed \equiv 1 \pmod{\phi(n)}$ (here $17 \times 2753 \equiv 1
\pmod{3120}$), a theorem built on Euler's theorem guarantees $(m^e)^d \bmod n = m$ for any $m$
coprime to $n$. Concretely, $65^{17} \bmod 3233 = 2790$, and $2790^{2753} \bmod 3233$ returns exactly
$65$ again, because the modular exponentiation with $d$ is mathematically constructed to be the
exact inverse operation of exponentiation with $e$ under modulus $n$.
\answerlabel{5.4.A}
\item The performance reason: RSA exponentiation cost grows with the size of the input, so signing a
short, fixed-size hash rather than an arbitrarily long message keeps the expensive private-key
operation cheap and constant-size regardless of message length. The structural-security reason: an
attacker who obtains RSA signatures on certain chosen raw messages can sometimes combine them
algebraically to forge a signature on a message never actually signed. Hashing first removes this
algebraic structure by feeding the signature operation an unpredictable, avalanche-effect digest
instead of the raw message.
\answerlabel{5.4.A}
\item HMAC's shared-secret key does not distinguish which of the two key-holders produced a given
tag, since either party could have computed it -- there is no third piece of evidence separating
Alice from Bob. Digital signatures use different keys for signing and verifying: only the actual
signer's private key can produce a valid signature. So a signature that verifies against Bob's
public key could only have come from Bob, resolving the ambiguity structurally rather than through
trust between the two parties.
\answerlabel{5.4.A}
\item Trust in a self-signed certificate depends entirely on whether its signing key is already
present in the verifier's trust store through some independent, out-of-band decision. A public
website's self-signed certificate is not signed by any CA in the browser's pre-installed trust
store, so the chain of trust cannot be validated and the browser rejects it. The internal network's
self-signed certificate was manually installed into that specific browser's trust store by an
administrator, which is exactly the out-of-band act of trust a root CA normally provides at
OS-vendor scale. So the identical cryptographic construction is accepted there and nowhere else.
\answerlabel{5.4.C}
\item Validation requires each certificate's stated issuer to match the subject of the certificate
that actually signed it, forming an unbroken chain of names as well as an unbroken chain of
signatures. If the leaf's \texttt{issuer} field does not match the intermediate's \texttt{subject}
field, the chain link itself is broken -- there is no defined next certificate to verify the leaf's
signature against, regardless of whether some other, unrelated certificate's signature happens to
be mathematically valid. A validator following the chain-of-trust logic returns a broken-chain-link
failure at exactly this point, and never reaches the question of cryptographic validity for the
mismatched link at all.
\answerlabel{5.4.C}
\item Key-size comparisons describe how hard it is to break a key through cryptanalysis or
brute-force; they say nothing about a scenario where the key has already been stolen outright, since
a stolen key requires no cryptanalysis at all -- the attacker simply has it. What determines the
scope of damage here is the chain of trust's structure: every certificate anywhere that chains up to
this root, through any number of intermediate CAs, is now falsifiable. An attacker holding
the root's private key can forge a valid signature for a certificate claiming to be any website at
all, and every browser with that root in its trust store will accept it without complaint until the
compromise is disclosed and the root is actively removed or revoked.
\answerlabel{5.4.B}
\end{answerkey}

\section*{Sources}
\begin{enumerate}[label=\lbrack\arabic*\rbrack, leftmargin=2.2em]
\item \srcitem{5.4}{1}National Institute of Standards and Technology, \emph{Recommendation for Key
Management: Part 1 -- General}, NIST Special Publication 800-57 Part 1 Revision 5, May 2020, Table~2
(``Comparable security strengths of symmetric block cipher and asymmetric-key algorithms''), p.~54.
\item \srcitem{5.4}{2}National Institute of Standards and Technology, \emph{Recommendation for
Transitioning the Use of Cryptographic Algorithms and Key Lengths}, NIST Special Publication 800-131A,
January 2011 -- setting a transition deadline of the end of 2013, after which 80-bit-strength
algorithms including 1024-bit RSA are disallowed for applying cryptographic protection.
\end{enumerate}
